\documentclass[10pt,a4paper]{amsart}
\usepackage{amsmath,amssymb,amsthm,graphicx}
% Retained preamble for the Stacks Project chapter snapshot in this directory.
% The chapter's own preamble.tex is a copy of the shared Stacks Project preamble; the block
% below is sliced MECHANICALLY, line for line, from it. Nothing was added, removed or
% redefined; the wrapper class, the xr-hyper/hyperref wiring and the externaldocument list
% are omitted because the assembler supplies the document class.
%
% lines taken from preamble.tex: 3, 5, 9, 11, 12, 156, 157, 158, ... (51 lines total)
%
% package coverage restored from the chapter's own preamble (2026-10-13): xy, xy-2cell

\usepackage[all]{xy}
\xyoption{2cell}
\UseAllTwocells

\newenvironment{reference}{\comment}{\endcomment}
\newenvironment{slogan}{\comment}{\endcomment}
\newenvironment{history}{\comment}{\endcomment}
\theoremstyle{plain}
\newtheorem{theorem}[subsection]{Theorem}
\newtheorem{proposition}[subsection]{Proposition}
\newtheorem{lemma}[subsection]{Lemma}
\theoremstyle{definition}
\newtheorem{definition}[subsection]{Definition}
\newtheorem{example}[subsection]{Example}
\newtheorem{exercise}[subsection]{Exercise}
\newtheorem{situation}[subsection]{Situation}
\theoremstyle{remark}
\newtheorem{remark}[subsection]{Remark}
\newtheorem{remarks}[subsection]{Remarks}
\numberwithin{equation}{subsection}
\def\lim{\mathop{\mathrm{lim}}\nolimits}
\def\colim{\mathop{\mathrm{colim}}\nolimits}
\def\Spec{\mathop{\mathrm{Spec}}}
\def\Hom{\mathop{\mathrm{Hom}}\nolimits}
\def\Ext{\mathop{\mathrm{Ext}}\nolimits}
\def\SheafHom{\mathop{\mathcal{H}\!\mathit{om}}\nolimits}
\def\SheafExt{\mathop{\mathcal{E}\!\mathit{xt}}\nolimits}
\def\Sch{\mathit{Sch}}
\def\Mor{\mathop{\mathrm{Mor}}\nolimits}
\def\Ob{\mathop{\mathrm{Ob}}\nolimits}
\def\Sh{\mathop{\mathit{Sh}}\nolimits}
\def\NL{\mathop{N\!L}\nolimits}
\def\CH{\mathop{\mathrm{CH}}\nolimits}
\def\proetale{{pro\text{-}\acute{e}tale}}
\def\etale{{\acute{e}tale}}
\def\QCoh{\mathit{QCoh}}
\def\Ker{\mathop{\mathrm{Ker}}}
\def\Im{\mathop{\mathrm{Im}}}
\def\Coker{\mathop{\mathrm{Coker}}}
\def\Coim{\mathop{\mathrm{Coim}}}
\DeclareMathSymbol{\boxtimes}{\mathbin}{AMSa}{"02}
\def\QCohstack{\mathcal{QC}\!\mathit{oh}}
\def\Cohstack{\mathcal{C}\!\mathit{oh}}
\def\Spacesstack{\mathcal{S}\!\mathit{paces}}
\def\Quotfunctor{\mathrm{Quot}}
\def\Hilbfunctor{\mathrm{Hilb}}
\def\Curvesstack{\mathcal{C}\!\mathit{urves}}
\def\Polarizedstack{\mathcal{P}\!\mathit{olarized}}
\def\Complexesstack{\mathcal{C}\!\mathit{omplexes}}
\def\Pic{\mathop{\mathrm{Pic}}\nolimits}
\def\Picardstack{\mathcal{P}\!\mathit{ic}}
\def\Picardfunctor{\mathrm{Pic}}
\def\Deformationcategory{\mathcal{D}\!\mathit{ef}}
\usepackage{cleveref}
\crefname{originallemma}{original lemma}{original lemmas}
\Crefname{originallemma}{Original lemma}{Original lemmas}
\makeatletter
\newlabel{example-universal-quasi-finite-syntomic}{{0FK8}{source}{}{sourcepointer-1}{}}
\newlabel{example-universal-quasi-finite-syntomic@cref}{{[originallemma][0][]0FK8}{[1][0][]source}{}{}{}}
\makeatother
\crefname{originallemma}{original lemma}{original lemmas}
\Crefname{originallemma}{Original lemma}{Original lemmas}
\makeatletter
\newlabel{lemma-syntomic-quasi-finite}{{0BWE}{source}{}{sourcepointer-2}{}}
\newlabel{lemma-syntomic-quasi-finite@cref}{{[originallemma][0][]0BWE}{[1][0][]source}{}{}{}}
\makeatother
\title{Stacks, Discriminants: lemma-locally-comes-from-universal}
\author{The Stacks Project Authors}
\date{stacks-project:discriminant.tex:local-snapshot-20261012; GNU FDL 1.2 or later}
\begin{document}
\maketitle
\noindent\textit{Editorial note.} Statement span is lines 2224-2242 of sources/stacks-discriminant/pinned-original-source.tex; proof span is lines 2244-2262, the complete author proof (17 body lines).\par\medskip
\noindent\textit{Editorial note.} Navigation stubs: 2 cited-but-not-included labels are supplied as non-load-bearing pointers with their pinned tags where the snapshot tag index has them.\par\medskip
\noindent\textit{Editorial note.} extra\_packages is empty on purpose: the retained preamble already loads xy with options and a second declaration clashes in the assembler wrapper.\par\medskip
\noindent\textit{Editorial note.} Text is retained verbatim, including the author's own typos and placeholders.\par\medskip
\noindent\textit{Editorial note.} Queue index 40 of queue\_discriminant.json (pb=591, stmt 2224-2242, proof 2244-2262).\par\medskip
\noindent\textit{Allowed standard prerequisites (editorial).}\begin{itemize}
\item Explicit prerequisite boundary. The cited inputs are prerequisites; the argument in the retained span is not replaced by any cited result.
\end{itemize}\medskip
\ifdefined\hypertarget\hypertarget{sourcepointer-1}{}\fi
\noindent\textit{Source reference.} Original lemma 0FK8: sources/stacks-discriminant/pinned-original-source.tex, source lines 2168--2185. Cited by the retained statement or proof but not part of the collected spans; supplied only as a navigation pointer. The argument does not depend on its content.\par\medskip
\ifdefined\hypertarget\hypertarget{sourcepointer-2}{}\fi
\noindent\textit{Source reference.} Original lemma 0BWE: sources/stacks-discriminant/pinned-original-source.tex, source lines 2005--2026. Cited by the retained statement or proof but not part of the collected spans; supplied only as a navigation pointer. The argument does not depend on its content.\par\medskip
\begin{lemma}
\label{lemma-locally-comes-from-universal}
Let $f : Y \to X$ be a morphism of schemes. If $f$ satisfies the equivalent
conditions of Lemma \ref{lemma-syntomic-quasi-finite} then for every
$y \in Y$ there exist $n, d$ and a commutative diagram
$$
\xymatrix{
Y \ar[d] &
V \ar[d] \ar[l] \ar[r] &
Y_{n, d} \ar[d] \\
X & U \ar[l] \ar[r] &
X_{n, d}
}
$$
where $U \subset X$ and $V \subset Y$ are open with $y \in V$,
where $Y_{n, d} \to X_{n, d}$
is as in Example \ref{example-universal-quasi-finite-syntomic}, and
where the square on the right hand side is cartesian.
\end{lemma}

\begin{proof}
By Lemma \ref{lemma-syntomic-quasi-finite}
we can choose $U$ and $V$ affine so that
$U = \Spec(R)$ and $V = \Spec(S)$ with
$S = R[y_1, \ldots, y_n]/(g_1, \ldots, g_n)$.
With notation as in Example \ref{example-universal-quasi-finite-syntomic}
if we pick $d$ large enough, then we can write each $g_i$ as
$g_i = \sum_{E \in T} g_{i, E}y^E$ with $g_{i, E} \in R$.
Then the map $A \to R$ sending $a_{i, E}$ to $g_{i, E}$
and the map $B \to S$ sending $x_i \to y_i$ give a cocartesian
diagram of rings
$$
\xymatrix{
S & B \ar[l] \\
R \ar[u] & A \ar[l] \ar[u]
}
$$
which proves the lemma.
\end{proof}

\end{document}
